\chapter*{M3: Processes, Threads, and Dispatch}{Section Author: Sahil Gupta}
\addcontentsline{toc}{chapter}{M3: Processes, Threads, and Dispatch}
\section{Architecture Overview}

Milestone 3 adds the ability to spawn and control user-space processes outside of \co{init}.
Our implementation spans two main layers: the spawn library (which constructs the
address space and dispatcher for a new process), and the process manager in
\co{usr/init} (which owns the global process table and has control over processes). The
relevant code can be found in the following files:

\begin{itemize}
    \itemsep-0.5em
    \item \co{lib/spawn/spawn.c}
    \item \co{lib/spawn/vspace.c}
    \item \co{lib/spawn/dispatcher.c}
    \item \co{include/spawn/spawn.h}
    \item \co{usr/init/proc_mgmt.c}
    \item \co{usr/init/proc_mgmt.h}
    \item \co{lib/aos/paging.c}
\end{itemize}

\section{Process Spawning}

\subsection{Virtual Address Space Setup}

The kernel requires that the vspace root capability be rooted in the child's own
cspace, not in \co{init}'s. This means the L0 page table must be created inside
the child's \co{PAGECN}, not allocated from \co{init}'s slots. At the same time,
\co{init} needs its own view of that same L0 to install page table entries during
loading, since Barrelfish's paging API only operates on capabilities in the
calling process's cspace. We resolve this by creating the L0 in \co{PAGECN}
slot 0 and copying a reference into \co{init}'s cspace. The child's
\co{spawninfo} keeps the child-rooted cap as \co{si->vspace} for the dispatcher
invocation, while \co{init}'s paging operations use the copied reference.

\medskip\noindent
The child's paging state is initialized with \co{paging_init_state_foreign},
which builds a fresh shadow page table tree from scratch. Unlike \co{init}'s own paging
state (which is set up once at boot and grows its slabs on demand), the foreign
variant pre-allocates fixed-size slab buffers for shadow page table nodes and
virtual region tracking. The resulting state starts with a single free virtual
region from \co{BASE_PAGE_SIZE} up to the maximum virtual address, leaving the
null page unmapped. \co{init}'s slot allocator is shared with the child's paging
state so that new intermediate page table nodes can be allocated in \co{init}'s
cspace on the child's behalf as segments are mapped in.

\subsection{ELF Segment Mapping}

Loading the ELF requires \co{init} to temporarily have write access to the
child's physical memory so it can copy segment data from the module image. For
each loadable segment, a fresh physical frame is allocated and mapped twice: once
into the child's address space at the ELF-specified virtual address, and once
into \co{init}'s address space as a writable scratch window for the copy. A
capability to the frame is also stored in the child's \co{segcn} so the child
retains ownership of the physical memory after \co{init}'s mapping is torn down.

\medskip\noindent
The \co{init}-side mappings cannot be freed immediately after each segment copy.
Destroying a frame cap while the corresponding PTE still exists would
desynchronize the kernel's page tables from our shadow tree, causing
\co{SYS_ERR_VNODE_SLOT_INUSE} on the next mapping operation. Instead,
\co{spawn_vspace_alloc_and_map} records each \co{init}-side mapping in a small array
on \co{spawninfo} as it goes. Once the entire ELF is loaded and the module is
unmapped, its counterpart \co{spawn_vspace_unmap_resources} tears all of them
down together, removing the PTEs and destroying the caps in a single pass.
Without this cleanup, \co{init} would leak a cspace slot and a virtual region on
every spawn, eventually exhausting its allocators.

\subsection{Dispatcher Setup}

The dispatcher frame is the shared control structure between the kernel and the
process. \co{init} maps it into both address spaces. The child-side address is stored
in \co{disp->udisp} so the kernel knows where to find it after the handoff, and
\co{init} uses its own mapping to fill in the initial fields before the process starts.
The dispatcher is initialized in the disabled state with the process name and PID
set, so the kernel's scheduler can identify it from the moment it is made
runnable.

\section{Process Management}
The process management library (\co{init/proc_mgmt.c}) is used to keep track of the state of spawned
processes on a core. To facilitate this, we keep a linked list process table which holds every process
spawned. Currently this table is not garbage collected upon process exit (because resource clean up broadly
is not implemented for the \co{spawn}/\co{proc_mgmt} libraries). This table allows us to implement functionality
such as \cl{$ ps} as well as process lookup by name. \co{proc_mgmt} acts as a small wrapper around
the \co{spawn} library which dispatches to the correct function while maintaining the process table. As discussed above,
the \co{spawn} library enables spawning processes but it also controls process lifetimes. Functionality such
as making a process runnable, suspending a process, terminating a process, and waiting on a process is implemented in \co{spawn.c}. The \co{proc_mgmt}
library is the wrapper and entry point for this functionality, invoked via RPC.

\section{Development Process}


